% GENERATED FILE. Edit canonical lessons, then run npm run generate:books.
% canonical-source-hash: fnv1a64:54d86e584558b818
% canonical-lessons: ML-C570-pharniccar, ML-C570-balkkani, ML-C570-vatakakkaran, ML-C570-alakkuka-laundry, ML-C570-minukkuka

\chapter{Furniture and Laundry}
\label{ch:ml-furniture-and-laundry}
\hlchaptermodality{570}

\begin{chapteropening}
\textbf{By the end of this chapter:} \emph{I can say furniture, a balcony, a tenant, to do the laundry and to polish.}

Complete the last lesson of chapter 570: I can say furniture, a balcony, a tenant, to do the laundry and to polish.
\end{chapteropening}

\section[\texorpdfstring{\ml{ഫർണിച്ചർ} (pharṇiccar)}{ഫർണിച്ചർ (pharṇiccar)}]{\ml{ഫർണിച്ചർ} (pharṇiccar) — furniture}

\label{lesson:ML-C570-pharniccar}



\begin{quote}
Before the new one: say the Malayalam for to advance, then the Malayalam for to fall behind.
\end{quote}

\subsection*{You'll want to know: \ml{ഫർണിച്ചർ}}
\textbf{\ml{ഫർണിച്ചർ}} — \emph{pharṇiccar} — \textquotedblleft{}furniture\textquotedblright{}.

\begin{quote}
\ml{പുതിയ} \ml{വീട്ടിലേക്ക്} \ml{ഞങ്ങൾ} \ml{കുറച്ച്} \ml{ഫർണിച്ചർ} \ml{വാങ്ങി}.
\end{quote}

\begin{quote}
\emph{putiya vīṭṭilēkkŭ ñaṅṅaḷ kuṟaccŭ pharṇiccar vāṅṅi.}
\end{quote}

\begin{quote}
We bought some furniture for the new house.
\end{quote}

Like the English word, it has no plural.

\begin{cousinweb}
\textbf{\ml{ഫർണിച്ചർ}} (\emph{pharṇiccar}) was borrowed from English \emph{furniture}.
\end{cousinweb}

\begin{grammarlens}[title={What you've built}]
One more word for describing what you lived and felt.
\end{grammarlens}

\subsection*{Your turn}
\begin{itemize}
\raggedright
  \item \emph{Say it:} \emph{pharṇiccar}
  \item \emph{Say it:} \emph{pharṇiccar}, once more
  \item \emph{Say it:} \emph{putiya vīṭṭilēkkŭ ñaṅṅaḷ kuṟaccŭ pharṇiccar vāṅṅi.} — We bought some furniture for the new house.
  \item \emph{Recall:} say the Malayalam for to advance, then the Malayalam for to fall behind, then say \emph{pharṇiccar} again
\end{itemize}

\begin{tcolorbox}[breakable,colback=teal!4,colframe=teal!35!black,title={Before you move on}]
What does \ml{ഫർണിച്ചർ} mean? (\textquotedblleft{}Furniture\textquotedblright{}.) Say it once more.
\end{tcolorbox}
\section[\texorpdfstring{\ml{ബാൽക്കണി} (bālkkaṇi)}{ബാൽക്കണി (bālkkaṇi)}]{\ml{ബാൽക്കണി} (bālkkaṇi) — a balcony}

\label{lesson:ML-C570-balkkani}



\begin{quote}
Before the new one: say the Malayalam for to fall behind, then the Malayalam for furniture.
\end{quote}

\subsection*{You'll want to know: \ml{ബാൽക്കണി}}
\textbf{\ml{ബാൽക്കണി}} — \emph{bālkkaṇi} — \textquotedblleft{}a balcony\textquotedblright{}.

\begin{quote}
\ml{വൈകുന്നേരം} \ml{ഞങ്ങൾ} \ml{ബാൽക്കണിയിൽ} \ml{ഇരുന്ന്} \ml{ചായ} \ml{കുടിച്ചു}.
\end{quote}

\begin{quote}
\emph{vaikunnēraṁ ñaṅṅaḷ bālkkaṇiyil irunnŭ cāya kuṭiccu.}
\end{quote}

\begin{quote}
In the evening we sat on the balcony and drank tea.
\end{quote}

In a flat; a house usually has a \ml{വരാന്ത} instead.

\begin{cousinweb}
\textbf{\ml{ബാൽക്കണി}} (\emph{bālkkaṇi}) was borrowed from English \emph{balcony}.
\end{cousinweb}

\begin{grammarlens}[title={What you've built}]
One more word for describing what you lived and felt.
\end{grammarlens}

\subsection*{Your turn}
\begin{itemize}
\raggedright
  \item \emph{Say it:} \emph{bālkkaṇi}
  \item \emph{Say it:} \emph{bālkkaṇi}, once more
  \item \emph{Say it:} \emph{vaikunnēraṁ ñaṅṅaḷ bālkkaṇiyil irunnŭ cāya kuṭiccu.} — In the evening we sat on the balcony and drank tea.
  \item \emph{Recall:} say the Malayalam for to fall behind, then the Malayalam for furniture, then say \emph{bālkkaṇi} again
\end{itemize}

\begin{tcolorbox}[breakable,colback=teal!4,colframe=teal!35!black,title={Before you move on}]
What does \ml{ബാൽക്കണി} mean? (\textquotedblleft{}A balcony\textquotedblright{}.) Say it once more.
\end{tcolorbox}
\section[\texorpdfstring{\ml{വാടകക്കാരൻ} (vāṭakakkāran)}{വാടകക്കാരൻ (vāṭakakkāran)}]{\ml{വാടകക്കാരൻ} (vāṭakakkāran) — a tenant}

\label{lesson:ML-C570-vatakakkaran}



\begin{quote}
Before the new one: say the Malayalam for furniture, then the Malayalam for a balcony.
\end{quote}

\subsection*{You'll want to know: \ml{വാടകക്കാരൻ}}
\textbf{\ml{വാടകക്കാരൻ}} — \emph{vāṭakakkāran} — \textquotedblleft{}a tenant\textquotedblright{}.

\begin{quote}
\ml{മുകളിലെ} \ml{നിലയിൽ} \ml{ഒരു} \ml{പുതിയ} \ml{വാടകക്കാരൻ} \ml{വന്നു}.
\end{quote}

\begin{quote}
\emph{mukaḷile nilayil oru putiya vāṭakakkāran vannu.}
\end{quote}

\begin{quote}
A new tenant has come to the upper floor.
\end{quote}

\ml{വാടക} rent + -\ml{ക്കാരൻ}, one who pays it.

\begin{cousinweb}
\textbf{\ml{വാടകക്കാരൻ}} (\emph{vāṭakakkāran}) is built from \ml{വാടക}, \textquotedblleft{}rent\textquotedblright{}, usually traced to Sanskrit \emph{bhāṭaka}, \textquotedblleft{}hire\textquotedblright{} (Monier-Williams), with the ending -\ml{ക്കാരൻ}.
\end{cousinweb}

\begin{grammarlens}[title={What you've built}]
One more word for describing what you lived and felt.
\end{grammarlens}

\subsection*{Your turn}
\begin{itemize}
\raggedright
  \item \emph{Say it:} \emph{vāṭakakkāran}
  \item \emph{Say it:} \emph{vāṭakakkāran}, once more
  \item \emph{Say it:} \emph{mukaḷile nilayil oru putiya vāṭakakkāran vannu.} — A new tenant has come to the upper floor.
  \item \emph{Recall:} say the Malayalam for furniture, then the Malayalam for a balcony, then say \emph{vāṭakakkāran} again
\end{itemize}

\begin{tcolorbox}[breakable,colback=teal!4,colframe=teal!35!black,title={Before you move on}]
What does \ml{വാടകക്കാരൻ} mean? (\textquotedblleft{}A tenant\textquotedblright{}.) Say it once more.
\end{tcolorbox}
\section[\texorpdfstring{\ml{അലക്കുക} (alakkuka)}{അലക്കുക (alakkuka)}]{\ml{അലക്കുക} (alakkuka) — to do the laundry}

\label{lesson:ML-C570-alakkuka-laundry}



\begin{quote}
Before the new one: say the Malayalam for a balcony, then the Malayalam for a tenant.
\end{quote}

\subsection*{You'll want to know: \ml{അലക്കുക}}
\textbf{\ml{അലക്കുക}} — \emph{alakkuka} — \textquotedblleft{}to do the laundry\textquotedblright{}.

\begin{quote}
\ml{ഞായറാഴ്ച} \ml{ഞാൻ} \ml{എല്ലാ} \ml{തുണികളും} \ml{അലക്കി}.
\end{quote}

\begin{quote}
\emph{ñāyaṟāḻca ñān ellā tuṇikaḷuṁ alakki.}
\end{quote}

\begin{quote}
On Sunday I washed all the clothes.
\end{quote}

Past: \ml{അലക്കി}, \emph{alakki}.

Washing clothes by beating them on a stone; the past is \ml{അലക്കി}.

\begin{cousinweb}
\textbf{\ml{അലക്കുക}} (\emph{alakkuka}) is native Dravidian: DEDR 246 gives it as to wash clothes by beating them (DEDR 246).
\end{cousinweb}

\begin{grammarlens}[title={What you've built}]
One more word for describing what you lived and felt.
\end{grammarlens}

\subsection*{Your turn}
\begin{itemize}
\raggedright
  \item \emph{Say it:} \emph{alakkuka}
  \item \emph{Say it:} \emph{alakkuka}, once more
  \item \emph{Say it:} \emph{ñāyaṟāḻca ñān ellā tuṇikaḷuṁ alakki.} — On Sunday I washed all the clothes.
  \item \emph{Recall:} say the Malayalam for a balcony, then the Malayalam for a tenant, then say \emph{alakkuka} again
\end{itemize}

\begin{tcolorbox}[breakable,colback=teal!4,colframe=teal!35!black,title={Before you move on}]
What does \ml{അലക്കുക} mean? (\textquotedblleft{}To do the laundry\textquotedblright{}.) Say it once more.
\end{tcolorbox}
\section[\texorpdfstring{\ml{മിനുക്കുക} (minukkuka)}{മിനുക്കുക (minukkuka)}]{\ml{മിനുക്കുക} (minukkuka) — to polish}

\label{lesson:ML-C570-minukkuka}



\begin{quote}
Before the new one: say the Malayalam for a tenant, then the Malayalam for to do the laundry.
\end{quote}

\subsection*{You'll want to know: \ml{മിനുക്കുക}}
\textbf{\ml{മിനുക്കുക}} — \emph{minukkuka} — \textquotedblleft{}to polish\textquotedblright{}.

\begin{quote}
\ml{അതിഥികൾ} \ml{വരുന്നതിന്} \ml{മുൻപ്} \ml{അമ്മ} \ml{പിച്ചള} \ml{വിളക്ക്} \ml{മിനുക്കി}.
\end{quote}

\begin{quote}
\emph{atithikaḷ varunnatinŭ munpŭ amma piccaḷa viḷakkŭ minukki.}
\end{quote}

\begin{quote}
Before the guests came, Mother polished the brass lamp.
\end{quote}

Past: \ml{മിനുക്കി}, \emph{minukki}.

Its relative \ml{മിനുസമുള്ള}, smooth, is already familiar.

\begin{cousinweb}
\textbf{\ml{മിനുക്കുക}} (\emph{minukkuka}) is native Dravidian: DEDR 4876 lists it with \ml{മിന്നുക}, \textquotedblleft{}to flash, to shine\textquotedblright{} (DEDR 4876).
\end{cousinweb}

\begin{grammarlens}[title={What you've built}]
One more word for describing what you lived and felt.
\end{grammarlens}

\subsection*{Your turn}
\begin{itemize}
\raggedright
  \item \emph{Say it:} \emph{minukkuka}
  \item \emph{Say it:} \emph{minukkuka}, once more
  \item \emph{Say it:} \emph{atithikaḷ varunnatinŭ munpŭ amma piccaḷa viḷakkŭ minukki.} — Before the guests came, Mother polished the brass lamp.
  \item \emph{Recall:} say the Malayalam for a tenant, then the Malayalam for to do the laundry, then say \emph{minukkuka} again
\end{itemize}

\begin{tcolorbox}[breakable,colback=teal!4,colframe=teal!35!black,title={Before you move on}]
What does \ml{മിനുക്കുക} mean? (\textquotedblleft{}To polish\textquotedblright{}.) Say it once more.
\end{tcolorbox}
